\begin{textsample}{NEG Dim 1 – ma – Score 4.00 – ma/ma\_em\_1.txt}  \label{ex:f1_neg_224}
Apresentação É com satisfação que apresentamos o Documento Curricular do Território Maranhense – Ensino Médio/Volume II.
Ele ilustra o resultado de um expressivo esforço coletivo que envolveu as equipes técnicas desta Secretaria de Educação, professores e especialistas.
Para além do atendimento das prerrogativas legais, na consideração da Lei nº 13.415, de 16 de fevereiro de 2017, que altera a Lei nº 9.394/96, de Diretrizes e Bases da Educação Nacional, e a Resolução CNE/CEB nº 3, de 21 de novembro de 2018, que atualiza as Diretrizes Curriculares Nacionais para o Ensino Médio, a partir desses marcos legais, este documento objetiva orientar as equipes escolares no desenvolvimento de suas práticas pedagógicas no âmbito das escolas.
Assim, este documento se estrutura a partir de uma importante reflexão sobre o ensino médio oferecido no estado do Maranhão, desde os seus desafios aos princípios educacionais que orientam o projeto escolar e as práticas pedagógicas.
Na sequência, trazemos um aprofundamento sobre as áreas de conhecimento para, então, introduzir o tema flexibilização curricular – enfaticamente presente no Novo Ensino Médio e toda a sua estruturação.

% matched lemmas: 
\end{textsample}
